\chapter{System Design and Architecture}
This chapter describes the overall system design and architecture of the FoodSense AI project, including the hardware pipeline, sensor compensation, AI model architectures, and the web dashboard.

\section{System Overview}
FoodSense AI is a multi-modal IoT + AI system for early food spoilage detection. The system consists of three main components:
\begin{itemize}
    \item Hardware pipeline (Arduino + sensors)
    \item AI engine (MLP + LSTM)
    \item Web-based dashboard
\end{itemize}

The hardware pipeline collects sensor data, performs on-device compensation, and transmits the data in CSV format. The AI engine processes this data to detect ammonia and predict remaining shelf life. The web dashboard provides a user interface for visualizing the results.

Figure \ref{fig:system_architecture} illustrates the overall architecture and data flow of the FoodSense AI system.

\begin{figure}[htb]
\centering
\begin{tikzpicture}[
    node distance=1.2cm,
    block/.style={rectangle, draw, fill=blue!5, text width=2.8cm, text centered, rounded corners, minimum height=1.0cm, font=\fontsize{8}{10}\selectfont},
    cloud/.style={rectangle, draw, fill=green!5, text width=2.2cm, text centered, rounded corners, minimum height=0.8cm, font=\fontsize{8}{10}\selectfont},
    arrow/.style={thick,->,>=stealth}
]
    % Nodes
    \node (food) [cloud] {Produce \\ \& Spoilage VOCs};
    \node (sensors) [block, right=of food] {\textbf{Sensor Assembly} \\ MQ135 Gas Sensor \\ DHT11 Temp/Hum \\ RGB-LDR Optical Cell};
    \node (mcu) [block, right=of sensors] {\textbf{Arduino Edge} \\ Compensation Math \\ Multi-wavelength Cycle};
    \node (stream) [block, below=of mcu, yshift=-0.5cm] {\textbf{Inference Server} \\ FastAPI Backend \\ PyTorch / TF Models};
    \node (dashboard) [block, left=of stream, xshift=-1.5cm] {\textbf{Cloud Dashboard} \\ Live Plots \& Metrics \\ Vision Spoilage Status};

    % Arrows
    \draw [arrow] (food) -- (sensors);
    \draw [arrow] (sensors) -- (mcu);
    \draw [arrow] (mcu) -- node[right, font=\fontsize{7}{9}\selectfont] {CSV Data} (stream);
    \draw [arrow] (stream) -- node[above, font=\fontsize{7}{9}\selectfont] {WebSockets / HTTP} (dashboard);
\end{tikzpicture}
\caption{Multimodal System Architecture and Data Flow}
\label{fig:system_architecture}
\end{figure}

\section{Sensor Data Pipeline}
The sensor data pipeline is responsible for acquiring, processing, and transmitting data from the sensors.

\subsection{On-Device Sensor Compensation}
A key novelty of FoodSense AI is that temperature and humidity compensation for the MQ135 sensor is performed directly on the Arduino microcontroller, rather than in post-processing. The compensation factor is calculated as:
\[
\text{corr} = 1 + 0.02(T - 20) + 0.01(H - 65)
\]
where \(T\) is the temperature in $^\circ$C and \(H\) is the humidity in \% RH. The corrected sensor resistance is then:
\[
\text{corrected\_Rs} = \frac{\text{Rs}}{\text{corr}}
\]
where Rs is the uncompensated sensor resistance.

\subsection{CSV Transmission}
The Arduino transmits the sensor data in CSV format over the serial port. Each line contains:
\[
\text{timestamp\_ms, temperature, humidity, nh3\_ppm, status}
\]
This format allows direct ingestion into Python using pandas.read\_csv, without the need for additional parsing middleware.

\section{AI Model Architectures}
The FoodSense AI system uses two complementary AI models for different tasks.

\subsection{MLP Gas Classifier}
The MLP classifier is used for point-in-time detection of ammonia, a primary marker of food spoilage. The model takes 128 features as input (from the UCI Gas Sensor Array Drift Dataset) and outputs a binary classification (ammonia present or not).

Key design choices for the MLP include:
\begin{itemize}
    \item Batch normalization after each dense layer to stabilize training
    \item BCEWithLogitsLoss with pos\_weight = 7.5 to handle class imbalance
    \item L2 regularization via weight\_decay = 1e-4 to prevent overfitting
\end{itemize}

\subsection{LSTM with Temporal Attention}
The LSTM with temporal attention is used for continuous shelf-life prediction. Spoilage is a process, not a snapshot, so the model needs to capture temporal dependencies across time. The model uses a 1-hour context window (last 30 readings) to predict the remaining shelf life in hours.

Key design choices for the LSTM include:
\begin{itemize}
    \item Huber loss instead of MSE to be less sensitive to outliers
    \item ReduceLROnPlateau scheduler to halve the learning rate when validation loss plateaus
    \item Gradient clipping (max\_norm = 1.0) to prevent exploding gradients
    \item Target normalization to [0,1] during training for better gradient flow
\end{itemize}

\section{Web Dashboard Design}
The web dashboard is built using HTML, CSS, and JavaScript with Tailwind CSS for styling. It provides a user-friendly interface for visualizing sensor data and AI model predictions.

The dashboard includes four main tabs:
\begin{itemize}
    \item Dashboard: Real-time sensor readings and AI predictions (simulated in the current version)
    \item Evaluation: Model performance metrics (Precision, Recall, F1-Score)
    \item Image Analysis: Upload or capture images for spoilage detection (planned Phase 2)
    \item Architecture: Advanced multiflow architecture diagram showing 8 components and branching tracks
\end{itemize}

\section{Novelty and Contributions}
FoodSense AI introduces several technical novelties compared to traditional systems:
\begin{itemize}
    \item \textbf{Multiflow Architecture}: A branched pipeline that processes gas and climate data through parallel MLP and LSTM tracks for simultaneous status and shelf-life inference.
    \item \textbf{On-Device Compensation}: Signal correction performed at the edge (microcontroller) to ensure data integrity before transmission.
    \item \textbf{Automated Cloud Sync}: A CI/CD pipeline using GitHub Actions that automatically synchronizes model updates and code to the live Hugging Face Space.
\end{itemize}

\section{Summary}
This chapter described the overall system design and architecture of the FoodSense AI project, including the sensor data pipeline, the multiflow AI engine, and the web-based cloud dashboard. The next chapter discusses the implementation details.
